\documentclass[11pt,oneside]{article}
\usepackage[a4paper,margin=27mm,headheight=15pt]{geometry}
\usepackage[T1]{fontenc}
\usepackage[utf8]{inputenc}
\IfFileExists{libertinus.sty}{\usepackage{libertinus}}{\usepackage{lmodern}}
\usepackage{microtype}
\usepackage{amsmath,amssymb,amsthm,mathtools,mathrsfs}
\usepackage{bm}
\usepackage{booktabs,longtable,array,tabularx}
\usepackage{graphicx}
\usepackage{pgfplots}
\usepackage{xcolor}
\usepackage{enumitem}
\usepackage{fancyhdr}
\usepackage{titlesec}
\usepackage{xurl}
\usepackage{aliascnt}
\usepackage{hyperref}
\usepackage[nameinlink,noabbrev]{cleveref}
\definecolor{linkblue}{RGB}{25,75,135}
\definecolor{shadegray}{RGB}{246,247,249}
\definecolor{rulegray}{RGB}{195,201,208}
\pgfplotsset{compat=1.18}
\hypersetup{
  colorlinks=true,
  linkcolor=linkblue,
  citecolor=linkblue,
  urlcolor=linkblue,
  pdftitle={Sharp Real-Axis Decay of the Fourier Transform of Rvachev's Up-Function},
  pdfsubject={Dyadic sinc products, Thue-Morse dynamics, Gelfond extremals, Lambert W asymptotics},
  pdfkeywords={Rvachev up-function, Fabius function, sinc product, Fourier decay, Thue-Morse, Gelfond exponent, Lambert W}
}
\pagestyle{fancy}
\fancyhf{}
\fancyhead[L]{\small Fourier decay of Rvachev's up-function}
\fancyhead[R]{\small\nouppercase{\leftmark}}
\fancyfoot[C]{\thepage}
\renewcommand{\headrulewidth}{0.4pt}
\titleformat{\section}{\Large\bfseries}{\thesection}{0.7em}{}
\titleformat{\subsection}{\large\bfseries}{\thesubsection}{0.7em}{}
\titleformat{\subsubsection}{\normalsize\bfseries}{\thesubsubsection}{0.7em}{}
\setcounter{secnumdepth}{3}
\setcounter{tocdepth}{2}
\setlist[itemize]{leftmargin=2em,itemsep=0.25em,topsep=0.4em}
\setlist[enumerate]{leftmargin=2.3em,itemsep=0.25em,topsep=0.4em}
\setlength{\emergencystretch}{2em}
\allowdisplaybreaks[2]

\newtheorem{theorem}{Theorem}[section]
\newaliascnt{proposition}{theorem}
\newtheorem{proposition}[proposition]{Proposition}
\aliascntresetthe{proposition}
\newaliascnt{lemma}{theorem}
\newtheorem{lemma}[lemma]{Lemma}
\aliascntresetthe{lemma}
\newaliascnt{corollary}{theorem}
\newtheorem{corollary}[corollary]{Corollary}
\aliascntresetthe{corollary}
\newaliascnt{conjecture}{theorem}
\newtheorem{conjecture}[conjecture]{Conjecture}
\aliascntresetthe{conjecture}
\theoremstyle{definition}
\newaliascnt{definition}{theorem}
\newtheorem{definition}[definition]{Definition}
\aliascntresetthe{definition}
\newaliascnt{example}{theorem}
\newtheorem{example}[example]{Example}
\aliascntresetthe{example}
\theoremstyle{remark}
\newaliascnt{remark}{theorem}
\newtheorem{remark}[remark]{Remark}
\aliascntresetthe{remark}
\newaliascnt{warning}{theorem}
\newtheorem{warning}[warning]{Warning}
\aliascntresetthe{warning}

\crefname{theorem}{theorem}{theorems}
\Crefname{theorem}{Theorem}{Theorems}
\crefname{proposition}{proposition}{propositions}
\Crefname{proposition}{Proposition}{Propositions}
\crefname{lemma}{lemma}{lemmas}
\Crefname{lemma}{Lemma}{Lemmas}
\crefname{corollary}{corollary}{corollaries}
\Crefname{corollary}{Corollary}{Corollaries}
\crefname{conjecture}{conjecture}{conjectures}
\Crefname{conjecture}{Conjecture}{Conjectures}
\crefname{definition}{definition}{definitions}
\Crefname{definition}{Definition}{Definitions}
\crefname{example}{example}{examples}
\Crefname{example}{Example}{Examples}
\crefname{remark}{remark}{remarks}
\Crefname{remark}{Remark}{Remarks}
\crefname{warning}{warning}{warnings}
\Crefname{warning}{Warning}{Warnings}

\newcommand{\R}{\mathbb R}
\newcommand{\C}{\mathbb C}
\newcommand{\N}{\mathbb N}
\newcommand{\Z}{\mathbb Z}
\newcommand{\sinc}{\operatorname{sinc}}
\newcommand{\Up}{\operatorname{up}}
\newcommand{\e}{\mathrm e}
\newcommand{\dd}{\,\mathrm d}
\newcommand{\bigO}{\mathcal O}
\newcommand{\smallo}{o}
\newcommand{\vtwo}{v_2}
\newcommand{\dist}{\operatorname{dist}}
\newcommand{\supp}{\operatorname{supp}}
\newcommand{\W}{W}
\newcommand{\PhiF}{\Phi}
\newcommand{\Qprod}{Q}
\newenvironment{boxedremark}{%
  \par\smallskip\noindent\begin{center}\begin{minipage}{0.94\textwidth}%
  \color{black}\hrule\vspace{0.6em}\small
}{%
  \vspace{0.6em}\hrule\end{minipage}\end{center}\smallskip
}

\title{\bfseries Sharp Real-Axis Decay of the Fourier Transform\\
of Rvachev's Up-Function\\[0.35em]
\large Dyadic dynamics, Gelfond extremals, Lambert $W$ boundary layers, and two-adic troughs}
\author{A rigorous analysis of the decay-rate conjecture\\
\small in Vladimir Reshetnikov's \texttt{ProveIt} research frontier}
\date{26 August 2026}

\begin{document}
\maketitle

\begin{abstract}
With the Fourier convention
\(
 \widehat f(\xi)=\int_{\R}f(t)\e^{-2\pi i t\xi}\dd t,
\)
the normalized Rvachev up-function has Fourier transform
\[
 \PhiF(\xi)=\prod_{h=0}^{\infty}
 \sinc\!\left(\frac{\pi\xi}{2^h}\right),
 \qquad \sinc z=\frac{\sin z}{z}.
\]
The product is rapidly decreasing, but it has a highly nonuniform real-axis
amplitude: it vanishes at every nonzero integer, its dyadic orbits are governed
by Birkhoff sums of \(\log|\sin \pi x|\), and its largest peaks in a dyadic
annulus arise from a boundary layer that must first escape the zero at the
left endpoint and then shadow the maximizing Thue--Morse orbit
\(\{1/3,2/3\}\).

This article gives a rigorous hierarchy of decay laws.  Universally,
\[
 \limsup_{x\to\infty}
 \frac{\log|\PhiF(x)|}{(\log x)^2}
 =-\frac{1}{2\log 2},
 \qquad
 \liminf_{x\to\infty}
 \frac{\log|\PhiF(x)|}{(\log x)^2}=-\infty.
\]
For almost every fixed dyadic ray \(x=2^k y\), \(1<y<2\), an exact second
term is obtained from Birkhoff's theorem.  More sharply, put
\[
 M_k=\max_{2^k\le x\le 2^{k+1}}|\PhiF(x)|,
 \quad \lambda=\log 2,
 \quad \beta=\log\frac{\sqrt3}{2},
 \quad c=\frac{\lambda}{\pi}\sqrt{\frac32},
\]
and let \(W\) be the principal Lambert function.  We prove
\begin{align*}
 \log M_k={}&-\frac{\lambda}{2}k(k+1)
 -(k+1)\log\pi+(k+1)\beta\\
 &-\frac{W(c(k+1))^2+2W(c(k+1))}{2\lambda}
 +\bigO\!\left((\log\log k)^2\right).
\end{align*}
Thus the envelope contains a second log-normal suppression in
\(\log\log x\); it is not a log-normal term multiplied merely by a fixed
power of \(\log x\).  An elementary algebraic subaction proves the needed
sharp Gelfond exponent.  Finally, if \(E_n\) is the unique interzero peak on
\((n,n+1)\), then for each fixed odd positive integer \(m\),
\[
 E_{m2^k}\sim
 \frac{H(1)|H(m)|}{\e(k+1)}(m2^k)^{-(k+1)},
 \qquad
 H(z)=\prod_{r=1}^{\infty}\sinc\!\left(\frac{\pi z}{2^r}\right).
\]
These two-adically exceptional peaks are exponentially smaller on the
log-normal scale than the largest peak in the same dyadic annulus.  The
results show why no positive smooth normalization can make the pointwise
oscillation have a nonzero constant amplitude, and they separate pointwise,
typical, extremal, and averaged notions of decay.
\end{abstract}

\tableofcontents
\newpage

\section{The question and the necessary distinction between decay notions}
\label{sec:introduction}

Rvachev's up-function is an even, compactly supported \(C^\infty\) bump.
Under the Fourier convention used above, its Fourier transform is the entire
function
\begin{equation}
 \PhiF(z)=\prod_{h=0}^{\infty}
 \sinc\!\left(\frac{\pi z}{2^h}\right).
 \label{eq:Phi-definition}
\end{equation}
This product already occurs in the classical construction by Jessen and
Wintner and in the later theory of the Fabius and Rvachev functions; a modern
self-contained treatment is given by Arias de Reyna \cite{arias2017}.  The
product representation, its cosine form, and its canonical zero product are
also summarized in the Fabius-function literature \cite{wikipediaFabius}.

The draft that motivated the present analysis \cite{proveitDraft} proposes a
single asymptotic of the schematic form
\[
 |\PhiF(x)|\asymp
 (\log x)^{-p}x^{-\log(2\pi)/\log 2}
 \exp\!\left(-\frac{(\log x)^2}{2\log2}\right),
\]
and suggests that a reciprocal normalization should leave an oscillation of
asymptotically constant amplitude.  The leading quadratic logarithm is in the
right family, but the statement cannot be correct as a pointwise equivalent:
\(\PhiF\) vanishes at every nonzero integer, and its nonzero peaks vary by an
additional log-normal factor even inside a single dyadic annulus.  Moreover,
the extremal envelope has a Lambert \(W\) correction of size
\(-\tfrac1{2\log2}(\log\log x)^2\), far larger than a fixed power of
\(\log x\).

The word ``amplitude'' can therefore refer to several inequivalent objects.
We shall distinguish the following.
\begin{enumerate}
 \item The universal pointwise logarithmic scale of \(\PhiF(x)\).
 \item The decay along a fixed dyadic ray \(x=2^k y\).
 \item The largest value
 \[
   M_k=\max_{2^k\le x\le 2^{k+1}}|\PhiF(x)|
 \]
 in a dyadic annulus.
 \item The unique local peak
 \[
   E_n=\max_{n\le x\le n+1}|\PhiF(x)|
 \]
 between two consecutive integer zeros.
 \item Integral or Cesaro averages, which depend on \(L^1\), rather than
 \(L^\infty\), information about Thue--Morse products.
\end{enumerate}
The first four are treated rigorously here.  The fifth is reduced to a
separate transfer-operator problem in \cref{sec:averages}.

\begin{boxedremark}
\textbf{Main verdict.}
The correct leading scale is log-normal:
\(-\log|\PhiF(x)|\) is of order \((\log x)^2\) at its large nonzero peaks.
There is, however, no single pointwise asymptotic.  Typical dyadic rays,
maximizing rays, dyadic-annulus maxima, and peaks following highly even
integers have different lower-order laws.  The sharp dyadic-annulus envelope
is the Lambert \(W\) formula in \cref{thm:annulus-main}; its expansion begins
with a second quadratic logarithm in \(\log\log x\).
\end{boxedremark}

Throughout the article,
\begin{equation}
 \lambda=\log2,
 \qquad
 \beta=\log\frac{\sqrt3}{2}=\frac12\log3-\log2<0,
 \qquad
 \log0=-\infty.
 \label{eq:constants-lambda-beta}
\end{equation}
All logarithms are natural.

\section{Product identities, zeros, and one peak per unit interval}
\label{sec:basic}

\subsection{Entire convergence and the scaling law}

Since
\(
 \sinc w=1-w^2/6+\bigO(w^4)
\)
near the origin, the product in \eqref{eq:Phi-definition} converges locally
uniformly on \(\C\).  It therefore defines an even entire function.  Shifting
the index gives the exact refinement law
\begin{equation}
 \PhiF(2z)=\sinc(2\pi z)\PhiF(z),
 \qquad
 \PhiF(z)=\sinc(\pi z)\PhiF(z/2).
 \label{eq:Phi-scaling}
\end{equation}
Repeated use of the first identity yields
\begin{equation}
 \PhiF(2^k y)=\PhiF(y)
 \prod_{j=1}^{k}\sinc(\pi2^j y).
 \label{eq:Phi-ray-product}
\end{equation}
This elementary formula is the bridge from the Fourier transform to the
doubling map.

Define the tail product
\begin{equation}
 H(z)=\prod_{r=1}^{\infty}
 \sinc\!\left(\frac{\pi z}{2^r}\right).
 \label{eq:H-definition}
\end{equation}
Then
\begin{equation}
 \PhiF(z)=\sinc(\pi z)H(z).
 \label{eq:Phi-H}
\end{equation}
The function \(H\) is entire.  It is positive on \([0,2)\), and
\begin{equation}
 h:=H(1)=-\PhiF'(1)
 =0.55377127588881009534\ldots.
 \label{eq:h-constant}
\end{equation}
The numerical value in \eqref{eq:h-constant} is included only for reference;
none of the proofs is numerical.

\subsection{Canonical zero product}

Euler's product
\[
 \frac{\sin\pi z}{\pi z}
 =\prod_{m=1}^{\infty}\left(1-\frac{z^2}{m^2}\right)
\]
can be inserted into every factor of \eqref{eq:Phi-definition}.  Reindexing
the locally uniformly convergent product gives
\begin{equation}
 \PhiF(z)=
 \prod_{m=1}^{\infty}
 \left(1-\frac{z^2}{m^2}\right)^{1+\vtwo(m)},
 \label{eq:canonical-product}
\end{equation}
where \(\vtwo(m)\) is the exponent of \(2\) in \(m\).  Indeed, a fixed
integer \(m\) occurs once for every \(h\) with \(2^h\mid m\), hence with
multiplicity \(1+\vtwo(m)\).  Formula \eqref{eq:canonical-product} appears in
Arias de Reyna's treatment \cite{arias2017}.

Consequently, the real zeros are precisely the nonzero integers, and the zero
at \(m\ne0\) has multiplicity
\begin{equation}
 \operatorname{ord}_{m}\PhiF=1+\vtwo(|m|).
 \label{eq:zero-multiplicity}
\end{equation}
The multiplicities are the first indication that powers of two are exceptional.

\subsection{Uniqueness of the interzero extremum}

The canonical product gives more than the zero set.  It implies that every
open unit interval contains exactly one critical point of \(\PhiF\), hence
exactly one peak of its absolute value.

\begin{theorem}[One peak between consecutive integer zeros]
\label{thm:unique-peak}
For every integer \(n\ge1\), there is a unique point
\(c_n\in(n,n+1)\) at which \(\PhiF'(c_n)=0\).  The function
\(|\PhiF|\) is strictly increasing on \((n,c_n)\) and strictly decreasing on
\((c_n,n+1)\).  In particular,
\[
 E_n:=|\PhiF(c_n)|
 =\max_{n\le x\le n+1}|\PhiF(x)|
\]
is well defined.
\end{theorem}

\begin{proof}
Away from the integer zeros, logarithmic differentiation of
\eqref{eq:canonical-product} is justified by locally uniform convergence and
gives
\begin{equation}
 \frac{\PhiF'(x)}{\PhiF(x)}
 =-2x\sum_{m=1}^{\infty}
 \frac{1+\vtwo(m)}{m^2-x^2}.
 \label{eq:logderivative}
\end{equation}
Put
\[
 S(x)=\sum_{m=1}^{\infty}
 \frac{1+\vtwo(m)}{m^2-x^2}.
\]
On \((n,n+1)\), termwise differentiation gives
\[
 S'(x)=2x\sum_{m=1}^{\infty}
 \frac{1+\vtwo(m)}{(m^2-x^2)^2}>0.
\]
The term \(m=n\) forces \(S(x)\to-\infty\) as \(x\downarrow n\), while
the term \(m=n+1\) forces \(S(x)\to+\infty\) as \(x\uparrow n+1\).
Thus \(S\) has one and only one zero.  By \eqref{eq:logderivative}, this is
the unique zero of \(\PhiF'/\PhiF\).  Its sign changes from positive to
negative, so \(\log|\PhiF|\), and hence \(|\PhiF|\), first increases and then
decreases.
\end{proof}

The theorem provides a stable numerical procedure: maximize the logarithm of
the product separately on each unit interval, then take the largest peak in
the desired dyadic annulus.  It also makes the contrast in
\cref{sec:two-adic-troughs} unambiguous: the exceptional small quantities
there are genuine local maxima, not samples accidentally taken near a zero.

\section{The universal leading logarithmic scale}
\label{sec:universal}

A first estimate already determines the coefficient of \((\log x)^2\).
It is sharp in the limsup sense but cannot be a pointwise equivalent because
of the zeros.

\begin{theorem}[Sharp universal leading coefficient]
\label{thm:universal-leading}
As \(x\to+\infty\),
\begin{equation}
 \limsup_{x\to\infty}
 \frac{\log|\PhiF(x)|}{(\log x)^2}
 =-\frac{1}{2\lambda},
 \qquad
 \liminf_{x\to\infty}
 \frac{\log|\PhiF(x)|}{(\log x)^2}=-\infty.
 \label{eq:limsup-leading}
\end{equation}
\end{theorem}

\begin{proof}
Let \(2^k\le x<2^{k+1}\).  For \(0\le h\le k\), use
\(|\sin u|\le1\), while for \(h>k\) use \(|\sinc u|\le1\).  Then
\begin{equation}
 |\PhiF(x)|
 \le \prod_{h=0}^{k}\frac{2^h}{\pi x}
 =\frac{2^{k(k+1)/2}}{(\pi x)^{k+1}}.
 \label{eq:coarse-upper}
\end{equation}
Because \(k=(\log x)/\lambda+\bigO(1)\), taking logarithms gives
\[
 \log|\PhiF(x)|
 \le-\frac{(\log x)^2}{2\lambda}+\bigO(\log x),
\]
which proves the upper bound for the limsup.

For the reverse inequality, take \(x_k=2^k/3\).  Iterating
\eqref{eq:Phi-scaling} and using
\(|\sin(\pi2^j/3)|=\sqrt3/2\) gives the exact identity
\begin{equation}
 |\PhiF(2^k/3)|
 =|\PhiF(1/3)|
 \left(\frac{3\sqrt3}{2\pi}\right)^k
 2^{-k(k+1)/2}.
 \label{eq:special-ray-one-third}
\end{equation}
Since \(\log(2^k/3)=k\lambda+\bigO(1)\), this sequence attains the
coefficient \(-1/(2\lambda)\).  Finally, \(\PhiF(n)=0\) for every positive
integer \(n\), so the liminf is \(-\infty\).
\end{proof}

\begin{corollary}[Between polynomial and stretched-exponential decay]
\label{cor:decay-class}
The function \(\PhiF\) is rapidly decreasing: for every \(N>0\),
\(|\PhiF(x)|=\bigO_N(x^{-N})\).  On the other hand, for every
\(\delta>0\), it is not true that
\(|\PhiF(x)|=\bigO(\e^{-x^\delta})\) along all sufficiently large real
\(x\); the sequence \eqref{eq:special-ray-one-third} is much larger than any
such stretched exponential.
\end{corollary}

\begin{proof}
The first assertion follows immediately from \eqref{eq:coarse-upper} by
choosing more factors than a prescribed power, and is also proved directly in
\cite{arias2017}.  The second follows by comparing
\((\log x)^2\) with \(x^\delta\) along \eqref{eq:special-ray-one-third}.
\end{proof}

The theorem answers only the leading logarithmic question.  The coefficient
of \(\log x\), and even the meaning of a second term, depends on how infinity
is approached.  The doubling dynamics makes that dependence explicit.

\section{Dyadic rays and Thue--Morse products}
\label{sec:rays}

\subsection{Exact dynamical decomposition}

Fix \(y=1+\theta\in(1,2)\).  Since \(2^j\) is even for \(j\ge1\),
\(
 \sin(\pi2^j y)=\sin(\pi2^j\theta).
\)
Taking logarithms in \eqref{eq:Phi-ray-product} gives the exact formula
\begin{equation}
\begin{split}
 \log|\PhiF(2^k y)|
 ={}&\log|\PhiF(y)|
 -\frac{\lambda}{2}k(k+1)
 -k\log(\pi y)\\
 &+\sum_{j=1}^{k}\log|\sin(\pi2^j\theta)|.
\end{split}
\label{eq:ray-log-exact}
\end{equation}
Thus every second-order statement is a statement about Birkhoff sums of the
integrable, unbounded potential
\begin{equation}
 \varphi(t)=\log|\sin\pi t|
 \label{eq:potential}
\end{equation}
under the doubling map \(T(t)=2t\pmod1\).

For later use, define the finite sine product
\begin{equation}
 Q_n(t)=\prod_{j=0}^{n-1}|\sin(\pi2^j t)|.
 \label{eq:Qn-definition}
\end{equation}
It is exactly the modulus of a Thue--Morse polynomial, up to the factor
\(2^{-n}\).  Indeed, if \(s_2(m)\) is the sum of the binary digits of \(m\),
then
\begin{equation}
 \prod_{j=0}^{n-1}(1-\e^{2\pi i2^j t})
 =\sum_{m=0}^{2^n-1}(-1)^{s_2(m)}\e^{2\pi i m t},
 \label{eq:TM-polynomial}
\end{equation}
and therefore
\begin{equation}
 2^nQ_n(t)=
 \left|\sum_{m=0}^{2^n-1}(-1)^{s_2(m)}\e^{2\pi i m t}\right|.
 \label{eq:Q-TM-modulus}
\end{equation}
The corresponding uniform exponent is the classical Gelfond exponent.  Its
special value needed here is proved elementarily in \cref{sec:gelfond}; the
general ergodic-optimization framework is developed by Fan, Schmeling, and
Shen \cite{fan2021}.

\subsection{A general ray law}

\begin{theorem}[Decay along a dyadic ray]
\label{thm:general-ray}
Let \(y=1+\theta\in(1,2)\), and suppose that the finite limit
\begin{equation}
 \alpha(\theta)=
 \lim_{k\to\infty}\frac1k
 \sum_{j=1}^{k}\log|\sin(\pi2^j\theta)|
 \label{eq:alpha-definition}
\end{equation}
exists.  Then, with \(x_k=2^k y\),
\begin{equation}
 \log|\PhiF(x_k)|
 =-\frac{(\log x_k)^2}{2\lambda}
 +\left(
   \frac{\alpha(\theta)}{\lambda}
   -\frac12-\frac{\log\pi}{\lambda}
  \right)\log x_k
 +\smallo(\log x_k).
 \label{eq:general-ray-asymptotic}
\end{equation}
\end{theorem}

\begin{proof}
Insert
\(
 \sum_{j=1}^{k}\log|\sin(\pi2^j\theta)|
 =k\alpha(\theta)+\smallo(k)
\)
into \eqref{eq:ray-log-exact}.  This gives
\[
 \log|\PhiF(2^k y)|
 =-\frac{\lambda}{2}k^2
 +k\left(\alpha(\theta)-\frac\lambda2-\log\pi-\log y\right)
 +\smallo(k).
\]
Substitute \(k=(\log x_k-\log y)/\lambda\).  The two terms containing
\((\log y)\log x_k\) cancel, leaving \eqref{eq:general-ray-asymptotic}.
\end{proof}

The cancellation explains why the coefficient in terms of the actual
frequency \(x_k\), rather than the lower endpoint \(2^k\), is independent of
the chosen mantissa \(y\).  The orbit average \(\alpha(\theta)\) is the only
remaining dynamical datum.

\subsection{The almost-everywhere ray law}

The potential \(\varphi\) in \eqref{eq:potential} belongs to \(L^1([0,1])\),
and
\begin{equation}
 \int_0^1\log|\sin\pi t|\dd t=-\log2=-\lambda.
 \label{eq:log-sine-integral}
\end{equation}
The doubling map is ergodic for Lebesgue measure, so Birkhoff's theorem applies.

\begin{corollary}[Typical dyadic ray]
\label{cor:typical-ray}
For Lebesgue-almost every \(y\in(1,2)\), as \(k\to\infty\),
\begin{equation}
 \log|\PhiF(2^k y)|
 =-\frac{\lambda}{2}k^2
 -k\left(\frac{3\lambda}{2}+\log(\pi y)\right)
 +\smallo(k).
 \label{eq:typical-k}
\end{equation}
Equivalently, for \(x_k=2^k y\),
\begin{equation}
 \log|\PhiF(x_k)|
 =-\frac{(\log x_k)^2}{2\lambda}
 -\left(\frac32+\frac{\log\pi}{\lambda}\right)\log x_k
 +\smallo(\log x_k).
 \label{eq:typical-x}
\end{equation}
The numerical value of the linear exponent is
\[
 \frac32+\frac{\log\pi}{\log2}
 =3.151496129472318\ldots.
\]
\end{corollary}

This almost-everywhere law is a natural interpretation of a ``typical''
amplitude.  It is not an envelope: exceptional orbits can have a much larger
Birkhoff average and therefore a larger Fourier amplitude.

\section{An elementary sharp Gelfond bound}
\label{sec:gelfond}

The maximal possible orbit average in \cref{thm:general-ray} is
\(\beta=\log(\sqrt3/2)\).  This is the \(q=2,c=1/2\) Gelfond exponent in the
terminology of \cite{fan2021}.  For the present problem it admits a short,
fully explicit proof.

\begin{theorem}[Uniform sine-product bound]
\label{thm:gelfond-bound}
For every \(t\in\R\) and every integer \(n\ge1\),
\begin{equation}
 Q_n(t)
 \le\sqrt{\frac53}\left(\frac{\sqrt3}{2}\right)^n.
 \label{eq:gelfond-Q-bound}
\end{equation}
The exponential rate is sharp: at \(t=1/3\) or \(2/3\),
\begin{equation}
 Q_n(t)=\left(\frac{\sqrt3}{2}\right)^n.
 \label{eq:gelfond-equality-rate}
\end{equation}
Consequently,
\begin{equation}
 \sup_{t\in\R}\limsup_{n\to\infty}\frac1n\log Q_n(t)
 =\beta.
 \label{eq:max-alpha-beta}
\end{equation}
\end{theorem}

\begin{proof}
Set
\[
 u_j=\sin^2(\pi2^j t).
\]
The double-angle formula gives the logistic recursion
\begin{equation}
 u_{j+1}=4u_j(1-u_j).
 \label{eq:logistic-recursion}
\end{equation}
Introduce the positive affine function
\[
 V(u)=1+\frac23u.
\]
A direct factorization yields
\begin{equation}
 3V(u)-4uV(4u(1-u))
 =\frac{(2u+1)(4u-3)^2}{3}\ge0
 \qquad(0\le u\le1).
 \label{eq:subaction-identity}
\end{equation}
Hence
\[
 2\sqrt{u_j}
 \le\sqrt3\left(\frac{V(u_j)}{V(u_{j+1})}\right)^{1/2}.
\]
Multiplying from \(j=0\) to \(n-1\) telescopes:
\[
 \prod_{j=0}^{n-1}2|\sin(\pi2^j t)|
 \le3^{n/2}\left(\frac{V(u_0)}{V(u_n)}\right)^{1/2}
 \le3^{n/2}\sqrt{\frac53}.
\]
Dividing by \(2^n\) proves \eqref{eq:gelfond-Q-bound}.

At \(t=1/3\), the doubling orbit alternates between \(1/3\) and \(2/3\), so
every \(u_j=3/4\).  Equality holds in the exponential rate and
\eqref{eq:gelfond-equality-rate} follows.
\end{proof}

\begin{remark}[A calibrated subaction]
Taking logarithms in the proof gives the pointwise inequality
\begin{equation}
 \log|\sin\pi t|
 \le\beta+\frac12\log V(\sin^2\pi t)
 -\frac12\log V(\sin^2 2\pi t).
 \label{eq:calibrated-subaction}
\end{equation}
It is an explicit calibrated subaction for the maximizing two-cycle.  The
factorization \eqref{eq:subaction-identity} also identifies the equality set:
apart from zero factors, equality occurs exactly when
\(\sin^2\pi t=3/4\), that is, on the orbit \(\{1/3,2/3\}\).
\end{remark}

\begin{corollary}[Maximizing dyadic rays]
\label{cor:maximizing-ray}
The largest finite value of \(\alpha(\theta)\) in
\eqref{eq:alpha-definition} is \(\beta\).  It is attained by
\(\theta=1/3,2/3\) and by their eventually periodic preimages.  Along such a
maximizing ray, \cref{thm:general-ray} has linear coefficient
\begin{equation}
 \frac{\beta}{\lambda}-\frac12-\frac{\log\pi}{\lambda}
 =-\frac{a}{\lambda},
 \qquad
 a:=\frac\lambda2+\log\pi-\beta
 =\log\!\left(\frac{2^{3/2}\pi}{\sqrt3}\right).
 \label{eq:a-constant}
\end{equation}
Numerically,
\[
 a=1.635144512355263\ldots,
 \qquad
 \frac{a}{\lambda}=2.359014879111741\ldots.
\]
\end{corollary}

For example, taking \(y=4/3\), so that \(\theta=1/3\), gives the exact
identity
\begin{equation}
 |\PhiF(2^k\tfrac43)|
 =|\PhiF(\tfrac43)|
 \left(\frac{\sqrt3}{2\pi(4/3)}\right)^k
 2^{-k(k+1)/2}.
 \label{eq:maximizing-ray-exact}
\end{equation}
The dyadic-annulus maximum is subtler.  A fixed maximizing mantissa such as
\(4/3\) is separated by a fixed proportion from the lower endpoint \(2^k\),
and the denominator \((1+\theta)^{k+1}\) then incurs an exponential loss in
\(k\).  The true maximum moves toward the endpoint while spending a slowly
growing number of doublings reaching the maximizing cycle.  That competition
is the source of Lambert \(W\).

\section{The sharp dyadic-annulus envelope}
\label{sec:annulus}

This section contains the main asymptotic result.  The proof has three
components:
\begin{enumerate}
 \item an exact reduction of \(\PhiF\) on \([2^k,2^{k+1}]\) to a weighted
 Thue--Morse sine product;
 \item an upper bound from \cref{thm:gelfond-bound} and a matching lower
 construction using dense preimages of the maximizing two-cycle;
 \item a two-variable optimization whose logarithmic coordinates tile the
 real line and reduce to the principal Lambert \(W\) function.
\end{enumerate}

\subsection{Exact annular reduction}

Write
\[
 x=2^k(1+\theta),
 \qquad 0\le\theta\le1,
 \qquad n=k+1.
\]
Using \eqref{eq:Phi-H} for the factor at scale \(2^k\) and then the first
\(k\) refinement factors gives
\begin{equation}
 |\PhiF(2^k(1+\theta))|
 =\frac{2^{-k(k+1)/2}}{\pi^n}
 H(1+\theta)(1+\theta)^{-n}Q_n(\theta).
 \label{eq:annulus-exact}
\end{equation}
Indeed, the numerator contains the factors
\(
 |\sin(\pi2^j\theta)|
\)
for \(0\le j\le k\), and the denominators multiply to
\(
 \pi^n(1+\theta)^n2^{k(k+1)/2}.
\)
Thus, if
\begin{equation}
 A_n=\sup_{0\le\theta\le1}
 H(1+\theta)(1+\theta)^{-n}Q_n(\theta),
 \label{eq:An-definition}
\end{equation}
then
\begin{equation}
 M_k=\pi^{-n}2^{-k(k+1)/2}A_n.
 \label{eq:Mk-An}
\end{equation}
All of the nontrivial lower-order behavior is contained in \(A_n\).

\subsection{Transition coordinates near the endpoint}

The factor \((1+\theta)^{-n}\) forces a maximizer toward \(\theta=0\), while
\(Q_n(0)=0\).  For \(0<\theta\le1/2\), choose an integer \(r\ge0\) and
\(s\in I:=[1/4,1/2]\) such that
\begin{equation}
 \theta=2^{-r}s.
 \label{eq:transition-coordinates}
\end{equation}
Endpoint ambiguities are harmless.  The integer \(r\) is the number of
initial doublings spent escaping the zero at \(0\).  If \(r<n\), then
\begin{equation}
 Q_n(2^{-r}s)=P_r(s)Q_{n-r}(s),
 \qquad
 P_r(s)=\prod_{h=1}^{r}\sin\!\left(\frac{\pi s}{2^h}\right).
 \label{eq:Q-transition-factorization}
\end{equation}
The small block has the uniform factorization
\begin{equation}
 P_r(s)
 =(\pi s)^r2^{-r(r+1)/2}H_r(s),
 \qquad
 H_r(s)=\prod_{h=1}^{r}
 \sinc\!\left(\frac{\pi s}{2^h}\right).
 \label{eq:small-block-factorization}
\end{equation}
On the compact interval \(I\), the quantities \(H_r(s)\) are bounded above
and below by positive absolute constants, uniformly in \(r\).  Hence
\begin{equation}
 \log P_r(s)
 =-\frac{\lambda}{2}r^2
 +r\left(\log(\pi s)-\frac\lambda2\right)
 +\bigO(1)
 \label{eq:small-block-log}
\end{equation}
uniformly for \(s\in I\).

The tail estimate from \cref{thm:gelfond-bound} gives
\begin{equation}
 Q_m(s)\le C_G\e^{m\beta},
 \qquad C_G=\sqrt{\frac53}.
 \label{eq:tail-upper}
\end{equation}
To match it from below while retaining freedom to choose \(s\), we use
preimages of the maximizing cycle.

\begin{lemma}[Dense near-maximizing preimages]
\label{lem:dense-preimages}
There is an absolute constant \(C>0\) with the following property.  For every
\(s_0\in I\) and every integer \(q\ge1\), one can choose
\(s_q\in I\) such that
\begin{equation}
 |s_q-s_0|\le2^{-q},
 \qquad
 T^q(s_q)\in\left\{\frac13,\frac23\right\},
 \label{eq:preimage-approximation}
\end{equation}
and, for every \(m\ge q\),
\begin{equation}
 Q_m(s_q)\ge\exp(m\beta-Cq^2).
 \label{eq:preimage-tail-lower}
\end{equation}
\end{lemma}

\begin{proof}
The \(q\)-fold preimages of \(1/3\) are
\[
 \frac{j+1/3}{2^q},
 \qquad 0\le j<2^q,
\]
and have spacing \(2^{-q}\).  Using these and, near an endpoint if needed,
the preimages of \(2/3\), choose \(s_q\in I\) satisfying
\eqref{eq:preimage-approximation}.

For \(0\le j<q\), the point \(T^j(s_q)\) is a nonintegral rational whose
reduced denominator divides \(3\cdot2^{q-j}\).  Therefore
\[
 \dist(T^j(s_q),\Z)\ge\frac{1}{3\cdot2^{q-j}}.
\]
The elementary inequality
\(
 |\sin\pi t|\ge2\dist(t,\Z)
\)
now gives
\[
 \prod_{j=0}^{q-1}|\sin(\pi T^j(s_q))|
 \ge\prod_{j=0}^{q-1}\frac{2}{3\cdot2^{q-j}}
 =\left(\frac23\right)^q2^{-q(q+1)/2}.
\]
After time \(q\), the orbit alternates between \(1/3\) and \(2/3\), so every
remaining factor equals \(\sqrt3/2=\e^\beta\).  Taking logarithms and
absorbing the linear terms into \(Cq^2\) proves
\eqref{eq:preimage-tail-lower}.
\end{proof}

The quadratic loss in \(q\) is more than sufficient: later we choose
\(q\asymp\log\log n\), while the main boundary-layer loss is of order
\((\log n)^2\).

\subsection{The finite-dimensional objective}

Combining \eqref{eq:small-block-log} with the tail rate \(\beta\), define
\begin{equation}
 \mathcal F_n(r,s)=
 -\frac{\lambda}{2}r^2
 +r\left(\log(\pi s)-\frac\lambda2-\beta\right)
 -n\log(1+s2^{-r}),
 \label{eq:Fn-exact}
\end{equation}
for integers \(r\ge0\) and \(s\in I\).

\begin{lemma}[Reduction to \(\mathcal F_n\)]
\label{lem:An-Fn-reduction}
As \(n\to\infty\),
\begin{equation}
 \log A_n
 =n\beta+
 \sup_{\substack{r\in\Z_{\ge0}\\s\in I}}
 \mathcal F_n(r,s)
 +\bigO\!\left((\log\log n)^2\right).
 \label{eq:An-Fn-reduction}
\end{equation}
\end{lemma}

\begin{proof}
We first record a benchmark.  Let
\(r_0=\lfloor(\log n)/\lambda\rfloor\), put
\(\theta_0=2^{-r_0}/3\), and use
\(Q_{n-r_0}(1/3)=\e^{(n-r_0)\beta}\) in
\eqref{eq:Q-transition-factorization}.  The uniform bounds on \(H\) and
\(H_{r_0}\) then give
\begin{equation}
 \log A_n\ge n\beta-C(\log n)^2
 \label{eq:An-benchmark}
\end{equation}
for an absolute constant \(C\).

For the upper bound, first note that \(0\le H(1+\theta)\le1\) on
\([0,1]\).  If \(\theta\ge1/2\), then
\[
 H(1+\theta)(1+\theta)^{-n}Q_n(\theta)
 \le C_G\exp\!\left(n\beta-n\log\frac32\right),
\]
which is exponentially smaller than the benchmark \eqref{eq:An-benchmark}.  If
\(0<\theta<1/2\), use \eqref{eq:transition-coordinates}.  When \(r<n\),
\eqref{eq:small-block-factorization} and \eqref{eq:tail-upper} give the
right-hand side of \eqref{eq:An-Fn-reduction}, with an \(\bigO(1)\) error.
When \(r\ge n\), all \(n\) sine factors are in their small-argument regime,
and
\[
 Q_n(\theta)\le(\pi\theta)^n2^{n(n-1)/2}
 \le\exp\!\left(-\frac\lambda2n^2+\bigO(n)\right),
\]
again negligible.  This proves the upper bound.

For the lower bound, choose \((r,s_0)\) within \(1\) of the supremum of
\(\mathcal F_n\).  We first record a quantitative localization.  With
\(r_0=\lfloor(\log n)/\lambda\rfloor\) and \(s=1/3\), one has
\(s2^{-r_0}=\bigO(n^{-1})\), and therefore
\[
 \mathcal F_n(r_0,1/3)\ge-C_1(\log n)^2
\]
for an absolute constant \(C_1\).  On the other hand, because
\(\log(\pi s)-\lambda/2-\beta\) is bounded on \(I\),
\[
 \mathcal F_n(r,s)\le-\frac\lambda2r^2+C_2r.
\]
Comparison with the preceding lower bound forces \(r=\bigO(\log n)\) at every
near-maximizer.  Since \(0\le s2^{-r}\le1/2\) and
\(\log(1+u)\ge u/2\) on this interval, we also have
\[
 \mathcal F_n(r,s)
 \le C_3r-\frac12ns2^{-r}.
\]
A second comparison with \(-C_1(\log n)^2\) now gives
\begin{equation}
 r=\bigO(\log n),
 \qquad
 n s_0 2^{-r}=\bigO((\log n)^2).
 \label{eq:Fn-localization-crude}
\end{equation}

Take
\[
 q=\left\lceil4\log_2(\log(n+\e^\e))\right\rceil
\]
after increasing it by a fixed constant if necessary.  For all sufficiently large
\(n\), the localization \eqref{eq:Fn-localization-crude} implies
\(n-r\ge q\).  Apply \cref{lem:dense-preimages} to obtain \(s_q\) with
\(|s_q-s_0|\le2^{-q}\).  The derivative of \(\mathcal F_n(r,s)\) with
respect to \(s\), throughout \(I\), is bounded by \(\bigO((\log n)^2)\) at
the localized values in \eqref{eq:Fn-localization-crude}.  Hence
\[
 \mathcal F_n(r,s_q)=\mathcal F_n(r,s_0)+\bigO(1).
\]
The factors \(H(1+2^{-r}s_q)\) and \(H_r(s_q)\) are bounded below by positive
constants.  Finally, \eqref{eq:preimage-tail-lower} supplies the tail with a
loss \(Cq^2=\bigO((\log\log n)^2)\).  This gives the lower bound in
\eqref{eq:An-Fn-reduction}.
\end{proof}

\subsection{Lambert \texorpdfstring{\(W\)}{W} optimization}

The remaining optimization has a useful exact structure.  Put
\begin{equation}
 C_0=\log\pi-\frac\lambda2-\beta
 =\log\!\left(\pi\sqrt{\frac23}\right),
 \qquad
 c=\lambda\e^{-C_0}
 =\frac{\lambda}{\pi}\sqrt{\frac32}.
 \label{eq:C0-c}
\end{equation}
Numerically,
\begin{equation}
 c=0.27022231973336536\ldots.
 \label{eq:c-numeric}
\end{equation}

\begin{lemma}[Lambert optimization]
\label{lem:Lambert-optimization}
Let \(w_n=W(cn)\), where \(W\) is the principal Lambert function.  Then
\begin{equation}
 \sup_{\substack{r\in\Z_{\ge0}\\s\in I}}
 \mathcal F_n(r,s)
 =-\frac{w_n^2+2w_n}{2\lambda}+\bigO(1).
 \label{eq:Lambert-optimization-result}
\end{equation}
\end{lemma}

\begin{proof}
First replace \(-n\log(1+s2^{-r})\) by \(-ns2^{-r}\), and denote the
resulting function by \(F_n(r,s)\).  At a maximizer of
\(\mathcal F_n\), the localization argument in
\eqref{eq:Fn-localization-crude} gives
\(ns2^{-r}=\bigO((\log n)^2)\).  Therefore
\[
 0\le n\bigl(s2^{-r}-\log(1+s2^{-r})\bigr)
 \le\frac n2s^22^{-2r}
 =\bigO\!\left(\frac{(\log n)^4}{n}\right)=\smallo(1).
\]
The same comparison applies to a maximizer of \(F_n\), so the two suprema
differ by \(\smallo(1)\).

Set
\[
 A=C_0+\log s.
\]
As \(s\) ranges over \(I=[1/4,1/2]\), the variable \(A\) ranges over an
interval \([A_-,A_+]\) of length exactly
\(
 A_+-A_-=\log2=\lambda.
\)
Now put
\begin{equation}
 z=\lambda r-A.
 \label{eq:z-variable}
\end{equation}
Since \(s=\e^{A-C_0}\), a direct calculation gives
\begin{equation}
 F_n(r,s)
 =-\frac{z^2}{2\lambda}
 -n\e^{-C_0}\e^{-z}
 +\frac{A^2}{2\lambda}.
 \label{eq:F-z-form}
\end{equation}
For successive integers \(r\), the intervals of possible \(z\) in
\eqref{eq:z-variable} have length \(\lambda\) and meet end to end.  Thus they
tile the real line, apart from an irrelevant lower truncation caused by
\(r\ge0\).  The last term in \eqref{eq:F-z-form} is uniformly bounded.
Consequently, the supremum differs by \(\bigO(1)\) from the maximum over real
\(z\) of
\[
 \phi_n(z)=-\frac{z^2}{2\lambda}-n\e^{-C_0}\e^{-z}.
\]
The critical-point equation is
\[
 -\frac z\lambda+n\e^{-C_0}\e^{-z}=0,
 \qquad\text{or equivalently}\qquad
 z\e^z=\lambda n\e^{-C_0}=cn.
\]
Hence the unique maximizer is \(z=w_n=W(cn)\).  At that point,
\(
 n\e^{-C_0}\e^{-w_n}=w_n/\lambda,
\)
so
\[
 \phi_n(w_n)
 =-\frac{w_n^2}{2\lambda}-\frac{w_n}{\lambda}
 =-\frac{w_n^2+2w_n}{2\lambda}.
\]
This proves \eqref{eq:Lambert-optimization-result}.
\end{proof}

The tiling argument is the reason the integer transition time \(r\) does not
create an \(\bigO(\log n)\) rounding error.  The continuous variable \(s\),
whose allowed interval has multiplicative width exactly two, compensates for
the lattice step in \(r\).  The remaining loss in
\cref{lem:An-Fn-reduction} comes from forcing the tail orbit to be genuinely
near-maximizing.

\subsection{Main envelope theorem}

\begin{theorem}[Sharp dyadic-annulus envelope]
\label{thm:annulus-main}
Let
\[
 M_k=\max_{2^k\le x\le2^{k+1}}|\PhiF(x)|,
 \qquad n=k+1,
 \qquad w_k=W(cn),
\]
where \(c\) is given by \eqref{eq:C0-c}.  Then, as \(k\to\infty\),
\begin{equation}
\boxed{
 \log M_k=
 -\frac{\lambda}{2}k(k+1)-n\log\pi+n\beta
 -\frac{w_k^2+2w_k}{2\lambda}
 +\bigO\!\left((\log\log k)^2\right).}
 \label{eq:annulus-main-centered}
\end{equation}
Equivalently, with \(a\) from \eqref{eq:a-constant},
\begin{equation}
 \log M_k=
 -\frac\lambda2k^2-ak
 -\frac{W(c(k+1))^2+2W(c(k+1))}{2\lambda}
 +\bigO\!\left((\log\log k)^2\right).
 \label{eq:annulus-main-compact}
\end{equation}
\end{theorem}

\begin{proof}
Combine \eqref{eq:Mk-An}, \cref{lem:An-Fn-reduction}, and
\cref{lem:Lambert-optimization}.  The constant difference between
\eqref{eq:annulus-main-centered} and \eqref{eq:annulus-main-compact} is
absorbed by the error term.
\end{proof}

The Lambert form is more accurate and more informative than a truncated
series.  Expanding
\begin{equation}
 W(cn)=\log n-\log\log n+\log c
 +\bigO\!\left(\frac{\log\log n}{\log n}\right)
 \label{eq:W-expansion}
\end{equation}
gives the following explicit hierarchy.

\begin{corollary}[Iterated-log expansion]
\label{cor:annulus-expanded}
As \(k\to\infty\),
\begin{equation}
\begin{split}
 \log M_k={}&-\frac\lambda2k^2-ak
 -\frac{(\log k)^2}{2\lambda}
 +\frac{\log k\,\log\log k}{\lambda}\\
 &-\frac{1+\log c}{\lambda}\log k
 +\bigO\!\left((\log\log k)^2\right).
\end{split}
\label{eq:annulus-expanded-k}
\end{equation}
\end{corollary}

In terms of the physical frequency scale, let \(R=2^k\),
\(L=\log R\), and
\begin{equation}
 d=\frac1\pi\sqrt{\frac32}.
 \label{eq:d-constant}
\end{equation}
Then \(c(k+1)=d(L+\lambda)\), and
\begin{equation}
\begin{split}
 \log\max_{R\le x\le2R}|\PhiF(x)|
 ={}&-\frac{L^2}{2\lambda}-\frac{a}{\lambda}L\\
 &-\frac{W(d(L+\lambda))^2+2W(d(L+\lambda))}{2\lambda}
 +\bigO\!\left((\log\log L)^2\right).
\end{split}
\label{eq:annulus-R-Lambert}
\end{equation}
Since \(W(dL)\sim\log L=\log\log R\), the first terms of the additional
boundary-layer correction are
\begin{equation}
 -\frac{(\log\log R)^2}{2\lambda}
 +\frac{\log\log R\,\log\log\log R}{\lambda}
 +\bigO(\log\log R).
 \label{eq:second-lognormal}
\end{equation}
This is the decisive correction missed by a normalization containing only a
fixed power of \(\log R\).

\subsection{Geometry of a near-maximizing point}

The proof also explains where the large peak lies.  At the smooth Lambert
optimum, the variable \(z\) is \(w_n\), and
\[
 n\theta=n s2^{-r}
 =n\e^{-C_0}\e^{-z}
 =\frac{w_n}{\lambda}.
\]
The preimage approximation in \cref{lem:dense-preimages} changes this only by
lower-order factors.  Thus one can construct near-maximizing points with
\begin{equation}
 \theta\asymp\frac{W(cn)}{n},
 \qquad
 r=\frac{W(cn)}{\lambda}+\bigO(1)
 =\log_2 n-\log_2\log n+\bigO(1).
 \label{eq:near-max-location}
\end{equation}
The frequency therefore has the form
\begin{equation}
 x=2^k\left(1+\Theta\!\left(\frac{\log k}{k}\right)\right).
 \label{eq:near-max-x}
\end{equation}
The orbit spends \(r\) steps in the small-angle region, then enters a compact
part of the circle and shadows an eventually periodic preimage of the
maximizing cycle.  The cost of the first phase is quadratic in \(r\); the
endpoint penalty is of order \(n2^{-r}\).  Balancing these two costs is exactly
the equation solved by Lambert \(W\).

\section{Two-adic troughs in the sequence of local peaks}
\label{sec:two-adic-troughs}

The dyadic-annulus maximum selects the best peak among exponentially many
unit intervals.  The individual peak sequence \(E_n\) is much less regular.
The most dramatic exceptions occur immediately after integers with large
two-adic valuation.

\subsection{Exact factorization near an odd multiple of a power of two}

Fix an odd positive integer \(m\), put
\[
 N=m2^k,
 \qquad x=N+t,
 \qquad 0\le t\le1,
\]
and define the finite product
\begin{equation}
 \PhiF_k(t)=\prod_{j=0}^{k}
 \sinc\!\left(\frac{\pi t}{2^j}\right).
 \label{eq:Phi-k-finite}
\end{equation}
For \(0\le j\le k\), the integer \(N/2^j\) contributes only a sign to the
sine.  Taking absolute values gives
\[
 \left|
 \sinc\!\left(\frac{\pi(N+t)}{2^j}\right)
 \right|
 =\frac{t}{N+t}
 \sinc\!\left(\frac{\pi t}{2^j}\right).
\]
The remaining factors are the tail \(H(m+t/2^k)\).  Hence the exact identity
\begin{equation}
 |\PhiF(N+t)|
 =\left(\frac{t}{N+t}\right)^{k+1}
 \PhiF_k(t)
 \left|H\!\left(m+\frac{t}{2^k}\right)\right|.
 \label{eq:near-m2k-factorization}
\end{equation}
Because \(m\) is odd, \(H(m)\ne0\).

As \(k\to\infty\), uniformly for \(0\le t\le1\),
\begin{equation}
 \PhiF_k(t)=\PhiF(t)\bigl(1+\bigO(4^{-k})\bigr),
 \qquad
 H\!\left(m+\frac{t}{2^k}\right)
 =H(m)\bigl(1+\bigO(2^{-k})\bigr),
 \label{eq:uniform-tail-approximations}
\end{equation}
The first estimate follows by writing \(\PhiF=\PhiF_k R_k\), where the
omitted tail \(R_k\) is positive on \([0,1]\) and satisfies
\(R_k=1+\bigO(4^{-k})\); thus the identity remains valid at \(t=1\), where
both sides vanish.  The second estimate follows from analyticity and
\(H(m)\ne0\).  Finally,
\begin{equation}
 \left(1+\frac{t}{N}\right)^{-(k+1)}=1+\smallo(1).
 \label{eq:N-denominator-uniform}
\end{equation}
Thus the peak problem reduces to maximizing \(t^p\PhiF(t)\), where
\(p=k+1\) tends to infinity.

\subsection{A boundary maximum lemma}

\begin{lemma}[Linear zero against a large power]
\label{lem:boundary-maximum}
Let \(f\) be continuous on \([0,1]\), positive on \([0,1)\), and suppose
\begin{equation}
 f(t)=h(1-t)+\bigO((1-t)^2)
 \qquad(t\uparrow1),
 \label{eq:linear-zero-assumption}
\end{equation}
with \(h>0\).  If
\[
 J_p=\max_{0\le t\le1}t^pf(t),
\]
then
\begin{equation}
 J_p\sim\frac{h}{\e p}.
 \label{eq:boundary-max-value}
\end{equation}
Moreover, if \(t_p\) is any maximizing point, then
\begin{equation}
 p(1-t_p)\longrightarrow1.
 \label{eq:boundary-max-location}
\end{equation}
\end{lemma}

\begin{proof}
Evaluation at \(t=1-1/p\), together with
\eqref{eq:linear-zero-assumption}, gives
\begin{equation}
 J_p\ge \left(1-\frac1p\right)^p
 \left(\frac hp+\bigO(p^{-2})\right)
 =\frac{h+\smallo(1)}{\e p}.
 \label{eq:boundary-lower-benchmark}
\end{equation}
For each fixed \(\varepsilon>0\), the part \(0\le t\le1-\varepsilon\)
contributes at most
\((1-\varepsilon)^p\|f\|_\infty=\smallo(p^{-1})\).  Hence every maximizing
point \(t_p\) tends to \(1\).

Put \(s=p(1-t)\).  Uniformly for \(s\) in a fixed compact subset of
\([0,\infty)\),
\begin{equation}
 p\,t^pf(t)\longrightarrow hs\e^{-s}.
 \label{eq:boundary-local-limit}
\end{equation}
It remains only to prevent the maximizing scaled variable
\(s_p=p(1-t_p)\) from escaping compact sets.  For \(t\) sufficiently close
to \(1\), \eqref{eq:linear-zero-assumption} gives
\(f(t)\le C(1-t)\), while
\(t^p\le\exp(-p(1-t))\).  Therefore
\begin{equation}
 p\,t^pf(t)\le C s\e^{-s}.
 \label{eq:boundary-tail-control}
\end{equation}
The right-hand side tends to zero both as \(s\downarrow0\) and as
\(s\to\infty\), whereas \eqref{eq:boundary-lower-benchmark} supplies a
positive limiting benchmark.  Thus \(s_p\) remains in a fixed compact
subset of \((0,\infty)\).  Every subsequential limit of \(s_p\) must, by
\eqref{eq:boundary-local-limit}, maximize \(hs\e^{-s}\).  This maximum is
unique and occurs at \(s=1\).  Hence \(s_p\to1\), and the maximum value is
asymptotic to \(h/(\e p)\).
\end{proof}

For \(f=\PhiF\), formula \eqref{eq:Phi-H} gives
\begin{equation}
 \PhiF(t)=h(1-t)+\bigO((1-t)^2),
 \qquad h=H(1),
 \label{eq:Phi-linear-at-one}
\end{equation}
because \((\sinc(\pi t))'|_{t=1}=-1\).

\begin{theorem}[Peaks following \(m2^k\)]
\label{thm:two-adic-peaks}
Fix an odd positive integer \(m\).  Let \(c_{m2^k}=m2^k+t_k\) be the unique
critical point in \((m2^k,m2^k+1)\).  Then
\begin{equation}
 (k+1)(1-t_k)\longrightarrow1,
 \label{eq:two-adic-peak-location}
\end{equation}
and
\begin{equation}
\boxed{
 E_{m2^k}\sim
 \frac{H(1)|H(m)|}{\e(k+1)}
 (m2^k)^{-(k+1)}.}
 \label{eq:two-adic-peak-main}
\end{equation}
In particular,
\begin{equation}
 E_{2^k}\sim
 \frac{H(1)^2}{\e(k+1)}2^{-k(k+1)}.
 \label{eq:powers-two-peak}
\end{equation}
\end{theorem}

\begin{proof}
Let \(p=k+1\).  From \eqref{eq:near-m2k-factorization}--
\eqref{eq:N-denominator-uniform}, uniformly in \(t\in[0,1]\),
\[
 |\PhiF(m2^k+t)|
 =(m2^k)^{-p}|H(m)|\,t^p\PhiF(t)(1+\smallo(1)).
\]
Apply \cref{lem:boundary-maximum} with \(f=\PhiF\) and
\eqref{eq:Phi-linear-at-one}.  The multiplicative error above tends to one
uniformly, so the benchmark and compactness argument in that lemma apply
unchanged to the perturbed maximization problem.  The uniqueness of the
maximizer follows from \cref{thm:unique-peak}, and the value and location
asymptotics follow.
\end{proof}

\begin{corollary}[Log-normal spread inside one annulus]
\label{cor:spread-annulus}
The peak immediately after \(2^k\) is exponentially smaller on the
log-normal scale than the largest peak in \([2^k,2^{k+1}]\):
\begin{equation}
 \log\frac{E_{2^k}}{M_k}
 =-\frac\lambda2k^2+\bigO(k).
 \label{eq:peak-spread}
\end{equation}
In particular, \(E_{2^k}/M_k\to0\) faster than every exponential in \(k\).
\end{corollary}

\begin{proof}
From \eqref{eq:powers-two-peak},
\(
 \log E_{2^k}=-\lambda k(k+1)-\log(k+1)+\bigO(1).
\)
Subtract \eqref{eq:annulus-main-compact}.
\end{proof}

This result rules out not only a pointwise constant-amplitude normalization,
but also any model in which all successive interzero peaks are assumed to
follow one smooth asymptotic curve.  The two-adic valuation remains visible
at arbitrarily large frequencies.

\section{What can and cannot be normalized}
\label{sec:normalization}

\subsection{Pointwise normalization}

No positive function \(g(x)\) can satisfy
\[
 \frac{|\PhiF(x)|}{g(x)}\longrightarrow C>0
 \qquad(x\to\infty)
\]
pointwise, because the quotient is zero at every positive integer.  Replacing
all points by local maxima does not restore a single regular law:
\cref{cor:spread-annulus} shows a log-normal spread among peaks whose
frequencies differ by a factor less than two.

What does exist is a hierarchy of precise normalizations:
\begin{center}
\begin{tabularx}{0.96\textwidth}{@{}>{\raggedright\arraybackslash}p{0.20\textwidth}X@{}}
\toprule
Object & Rigorous decay law \\
\midrule
Universal limsup &
\(\log|\PhiF(x)|\sim_{\limsup}-(\log x)^2/(2\log2)\).\\[0.4em]
Typical dyadic ray &
\(\log|\PhiF(x)|=-(\log x)^2/(2\log2)
 -(3/2+\log\pi/\log2)\log x+\smallo(\log x)\).\\[0.4em]
Maximizing fixed ray &
The same quadratic term, with linear coefficient
\(-a/\log2\), attained by the orbit \(\{1/3,2/3\}\).\\[0.4em]
Dyadic-annulus maximum &
The Lambert \(W\) formula \eqref{eq:annulus-main-centered}, including a
second log-normal correction in \(\log\log x\).\\[0.4em]
Peak after \(m2^k\) &
The much smaller two-adic law \eqref{eq:two-adic-peak-main}.\\
\bottomrule
\end{tabularx}
\end{center}

\subsection{Why the draft's proposed second term is not valid}

The draft \cite{proveitDraft} writes, in effect,
\begin{equation}
 \log|\PhiF(x)|
 \stackrel{?}{=}
 -\frac{(\log x)^2}{2\lambda}
 -\frac{\log(2\pi)}{\lambda}\log x
 -\frac12\log\log x+\bigO(1).
 \label{eq:draft-proposed}
\end{equation}
There are several independent obstructions.
\begin{enumerate}
 \item At the integer zeros, the left side is \(-\infty\).
 \item On almost every dyadic ray, the true linear coefficient is
 \(-3/2-\log\pi/\lambda\), not \(-\log(2\pi)/\lambda\); compare
 \eqref{eq:typical-x}.
 \item On maximizing rays the linear coefficient is again different;
 compare \eqref{eq:a-constant}.
 \item For the largest value in a dyadic annulus, the correction after the
 \(k^2\) and \(k\) terms is
 \[
  -\frac{(\log k)^2}{2\lambda}
  +\frac{\log k\log\log k}{\lambda}+\cdots,
 \]
 or, in the original frequency variable, the second log-normal term
 \eqref{eq:second-lognormal}.  It cannot be represented by
 \(-p\log\log x\) with constant \(p\).
 \item The peak sequence has the exceptional subsequence
 \eqref{eq:two-adic-peak-main}, whose leading quadratic coefficient in
 \(k\) is twice as negative as that of \(M_k\).
\end{enumerate}
The Fourier-series identity for \(\log|\sin|\) is useful only after the orbit
structure and its singularities are controlled.  Replacing the dyadic orbit
by its mean without specifying whether the problem is typical, extremal, or
averaged loses precisely the information that determines the lower-order
terms.

\section{The original Cesaro-average question is a different problem}
\label{sec:averages}

The Mathematics Stack Exchange question that motivated the draft asks for a
positive function \(g\) such that a Cesaro mean of \(|\PhiF|/g\) is bounded or
convergent \cite{reshetnikovMSEdecay}.  The present \(L^\infty\) analysis does
not by itself determine that normalization.

Indeed, integrating \eqref{eq:annulus-exact} gives
\begin{equation}
\begin{split}
 \int_{2^k}^{2^{k+1}}|\PhiF(x)|\dd x
 ={}&2^k\pi^{-n}2^{-k(k+1)/2}\,I_n,\\
 I_n={}&\int_0^1
 H(1+\theta)(1+\theta)^{-n}Q_n(\theta)\dd\theta.
\end{split}
\label{eq:annulus-L1-reduction}
\end{equation}
The quantity \(I_n\) is a weighted \(L^1\) norm of a Thue--Morse product.
By \eqref{eq:Q-TM-modulus}, the unweighted version is proportional to the
\(L^1\) norm of the length-\(2^n\) Thue--Morse polynomial, a separate and
subtle object discussed in the harmonic-analysis literature
\cite{queffelec2018}.  Its asymptotics are governed by pressure or
transfer-operator information, not solely by the Gelfond \(L^\infty\)
exponent.

Thus the sharp envelope derived here supplies an upper constraint and
identifies the rare configurations producing the largest spikes, but a
Cesaro normalizer requires additional information about the measure and
width of near-maximizing orbit sets.  Conflating these two questions is one
source of the incorrect pointwise claim in \eqref{eq:draft-proposed}.

\section{Numerical verification}
\label{sec:numerics}

The proofs above are analytic.  Numerical computation is nevertheless useful
for checking normalizations and for showing how early the Lambert formula
becomes accurate.

For each integer \(j\), \cref{thm:unique-peak} permits a one-dimensional
maximization of
\[
 \sum_{h=0}^{H}\log\left|
 \sinc\!\left(\frac{\pi x}{2^h}\right)\right|
\]
on \((j,j+1)\).  A tail bound from
\(
 \log\sinc u=-u^2/6+\bigO(u^4)
\)
controls the truncation.  Taking the largest of these unit-interval peaks
produces \(M_k\).  The table compares the computed value with the centered
Lambert approximation
\begin{equation}
 \mathcal A_k=
 -\frac\lambda2k(k+1)-(k+1)\log\pi+(k+1)\beta
 -\frac{W(c(k+1))^2+2W(c(k+1))}{2\lambda}.
 \label{eq:centered-numerical-approx}
\end{equation}

\begin{center}
\begin{tabular}{@{}rrrrr@{}}
\toprule
\(k\) & maximizing \(x/2^k\) & \(\log M_k\) & \(\mathcal A_k\) & residual \\
\midrule
3  & 1.177288 & -11.251778 & -10.428644 & -0.823133 \\
4  & 1.162272 & -15.485667 & -14.695388 & -0.790279 \\
5  & 1.168960 & -20.390823 & -19.639196 & -0.751626 \\
6  & 1.165501 & -25.991031 & -25.262965 & -0.728066 \\
7  & 1.167235 & -32.287929 & -31.568814 & -0.719116 \\
8  & 1.166376 & -39.274011 & -38.558353 & -0.715658 \\
9  & 1.166810 & -46.956338 & -46.232839 & -0.723499 \\
10 & 1.166594 & -55.329708 & -54.593276 & -0.736432 \\
11 & 1.166703 & -64.397553 & -63.640479 & -0.757074 \\
12 & 1.166649 & -74.157742 & -73.375125 & -0.782618 \\
\bottomrule
\end{tabular}
\end{center}

\begin{figure}[ht]
\centering
\begin{tikzpicture}
\begin{axis}[
 width=0.88\textwidth,
 height=0.48\textwidth,
 xlabel={dyadic index \(k\)},
 ylabel={logarithmic amplitude},
 xmin=3,xmax=12,
 grid=major,
 legend style={at={(0.03,0.03)},anchor=south west,draw=none,fill=none},
]
\addplot+[mark=*] coordinates {
 (3,-11.251778) (4,-15.485667) (5,-20.390823) (6,-25.991031)
 (7,-32.287929) (8,-39.274011) (9,-46.956338) (10,-55.329708)
 (11,-64.397553) (12,-74.157742)
};
\addlegendentry{computed \(\log M_k\)}
\addplot+[mark=square*] coordinates {
 (3,-10.428644) (4,-14.695388) (5,-19.639196) (6,-25.262965)
 (7,-31.568814) (8,-38.558353) (9,-46.232839) (10,-54.593276)
 (11,-63.640479) (12,-73.375125)
};
\addlegendentry{Lambert approximation \(\mathcal A_k\)}
\end{axis}
\end{tikzpicture}
\caption{The centered Lambert formula tracks the annular maximum with a
small bounded-size residual even at modest values of \(k\).  The theorem only
claims an \(\bigO((\log\log k)^2)\) remainder.}
\label{fig:numerical-comparison}
\end{figure}

For the displayed range, the maximizing ratio is close to \(7/6\).  This
corresponds to the simple preperiodic choice \(\theta=1/6\), which reaches the
maximizing cycle after one doubling.  At much larger \(k\), the optimal
preperiod grows according to \eqref{eq:near-max-location}; the convergence is
necessarily slow because it is only logarithmic in \(k\).

\section{Conclusions and a formalization roadmap}
\label{sec:conclusion}

The real-axis Fourier decay of the up-function is not described by one smooth
amplitude curve.  It has a clean but layered structure.

\begin{enumerate}
 \item The coefficient of \((\log x)^2\) is universally sharp in the limsup
 sense and equals \(-1/(2\log2)\).
 \item Fixed dyadic rays are controlled by Birkhoff averages of
 \(\log|\sin\pi x|\).  Almost every ray has average \(-\log2\), whereas the
 largest possible average is \(\log(\sqrt3/2)\).
 \item The sharp Gelfond rate follows from the elementary polynomial identity
 \eqref{eq:subaction-identity}; no computer-assisted ergodic optimization is
 needed in this special case.
 \item The largest peak in a dyadic annulus is a boundary-layer phenomenon.
 Its exact asymptotic through an \(\bigO((\log\log k)^2)\) remainder is the
 Lambert formula \eqref{eq:annulus-main-centered}.
 \item The first correction beyond the main log-normal envelope is itself
 log-normal in \(\log x\), namely quadratic in \(\log\log x\).  A fixed
 factor \((\log x)^{-p}\) cannot capture it.
 \item Individual local peaks retain strong two-adic irregularity.  The peaks
 after \(m2^k\) obey \eqref{eq:two-adic-peak-main} and are vastly smaller than
 the annular maximum.
\end{enumerate}

The proof was organized to be formalization-friendly.  A Lean development can
be divided into the following relatively independent blocks:
\begin{enumerate}
 \item locally uniform convergence, the scaling law, and the canonical zero
 product;
 \item monotonicity of the logarithmic-derivative sum and uniqueness of the
 interzero peak;
 \item the exact dyadic-ray and annular decompositions;
 \item the algebraic identity \eqref{eq:subaction-identity} and its telescoping
 consequence;
 \item the rational preimage construction in \cref{lem:dense-preimages};
 \item elementary localization estimates and the Lambert optimization in
 \cref{lem:Lambert-optimization};
 \item uniform convergence in \eqref{eq:uniform-tail-approximations} and the
 boundary maximum lemma.
\end{enumerate}
The only general ergodic input needed for the almost-everywhere statement is
Birkhoff's theorem.  The main extremal envelope and the two-adic peak theorem
are otherwise elementary consequences of the product itself.

\appendix

\section{Expanded forms and constants}
\label{app:constants}

For convenience, the constants occurring in the main theorem are
\begin{align*}
 \lambda&=\log2=0.6931471805599453\ldots,\\
 \beta&=\log(\sqrt3/2)=-0.1438410362258904\ldots,\\
 a&=\frac\lambda2+\log\pi-\beta
   =1.6351445123552633\ldots,\\
 c&=\frac{\lambda}{\pi}\sqrt{\frac32}
   =0.27022231973336536\ldots,\\
 d&=\frac1\pi\sqrt{\frac32}
   =0.3898484006168381\ldots,\\
 H(1)&=0.55377127588881009534\ldots.
\end{align*}

Let \(L_1=\log k\) and \(L_2=\log\log k\).  From the standard expansion of
Lambert \(W\) \cite{corless1996},
\[
 W(ck)=L_1-L_2+\log c
 +\frac{L_2-\log c}{L_1}
 +\bigO\!\left(\frac{L_2^2}{L_1^2}\right).
\]
Inserting only the terms justified above the remainder in
\cref{thm:annulus-main} gives \eqref{eq:annulus-expanded-k}.  Further formal
terms can be written, but the present proof's orbit-shadowing error is
\(\bigO(L_2^2)\), so coefficients at that order are not claimed.

\section{A compact proof of rapid decrease}
\label{app:rapid}

For completeness, fix \(N\ge1\).  From the first \(N\) factors of the product,
\[
 |x|^N|\PhiF(x)|
 \le |x|^N\prod_{h=0}^{N-1}
 \frac{2^h}{\pi|x|}
 =\frac{2^{N(N-1)/2}}{\pi^N}.
\]
Thus \(\PhiF\) is rapidly decreasing.  Differentiating the product and using
the same estimate, together with the geometric factors introduced by each
derivative, proves that every derivative of \(\PhiF\) is rapidly decreasing;
see \cite{arias2017} for the complete Schwartz-space argument.

\begin{thebibliography}{99}

\bibitem{arias2017}
J.~Arias de Reyna,
\newblock \emph{An infinitely differentiable function with compact support:
Definition and properties},
\newblock English translation and expanded version, arXiv:1702.05442, 2017.
\newblock \url{https://arxiv.org/abs/1702.05442}.

\bibitem{corless1996}
R.~M.~Corless, G.~H.~Gonnet, D.~E.~G.~Hare, D.~J.~Jeffrey, and D.~E.~Knuth,
\newblock On the Lambert \(W\) function,
\newblock \emph{Advances in Computational Mathematics} \textbf{5} (1996),
329--359.
\newblock \url{https://doi.org/10.1007/BF02124750}.

\bibitem{fan2021}
A.~Fan, J.~Schmeling, and W.~Shen,
\newblock \(L^\infty\)-estimation of generalized Thue--Morse trigonometric
polynomials and ergodic maximization,
\newblock \emph{Discrete and Continuous Dynamical Systems} \textbf{41}
(2021), no.~1, 297--327.
\newblock \url{https://doi.org/10.3934/dcds.2020363}.

\bibitem{gelfond1968}
A.~O.~Gelfond,
\newblock Sur les nombres qui ont des proprietes additives et
multiplicatives donnees,
\newblock \emph{Acta Arithmetica} \textbf{13} (1967/68), 259--265.
\newblock \url{https://doi.org/10.4064/aa-13-3-259-265}.

\bibitem{queffelec2018}
M.~Queffelec,
\newblock Questions around the Thue--Morse sequence,
\newblock \emph{Uniform Distribution Theory} \textbf{13} (2018), no.~1,
1--25.
\newblock \url{https://doi.org/10.1515/udt-2018-0001}.

\bibitem{reshetnikovMSEdecay}
V.~Reshetnikov,
\newblock Asymptotic decay rate of an infinite product of sinc functions,
\newblock Mathematics Stack Exchange, question 2745497, 20 April 2018.
\newblock \url{https://math.stackexchange.com/q/2745497}.

\bibitem{proveitDraft}
V.~Reshetnikov,
\newblock \emph{Decay rate conjecture},
\newblock non-formalized research-frontier draft in the \texttt{ProveIt}
repository, consulted 26 August 2026.
\newblock \url{https://github.com/VladimirReshetnikov/ProveIt/blob/main/Analysis/FabiusFunction/docs/non-formalized-research-frontiers/drafts/Decay%20rate%20conjecture/Decay%20rate%20conjecture.tex}.

\bibitem{wikipediaFabius}
Wikipedia contributors,
\newblock Fabius function,
\newblock consulted 26 August 2026.
\newblock \url{https://en.wikipedia.org/wiki/Fabius_function}.

\bibitem{walters1982}
P.~Walters,
\newblock \emph{An Introduction to Ergodic Theory},
\newblock Graduate Texts in Mathematics 79, Springer, 1982.

\end{thebibliography}

\end{document}
